\section{Results}
\subsection{Simulator VQE Accuracy}
Table \ref{table:distributed_results} presents the VQE ground state energy estimates within each of the 
benchmark molecules of $H_2, LiH, BeH_2, H_{2}O$, and compares it against the results of both the FCI reference and the serial baseline.

\begin{table}[htbp]
\centering
\renewcommand{\arraystretch}{1.2}
\caption{Distributed stack VQE results (P=2 MPI ranks, statevector simulator, CPU only, seed=42). The best physical energy is the lowest energy observed at/above the molecules FCI reference.}
\label{table:distributed_results}
\resizebox{\textwidth}{!}{%
\begin{tabular}{|l|c|c|c|r|r|r|c|r|}
\hline
\textbf{Molecule} & \textbf{Qubits} & \textbf{Pauli} & \textbf{Params} & \textbf{Energy ($E_h$)} & \textbf{FCI ($E_h$)} & \textbf{Error ($E_h$)} & \textbf{Chem. Acc.?} & \textbf{Iters} \\ \hline
H\textsubscript{2} & 4  & 15   & 24  & $-1.1349$ & $-1.1373$ & $0.0024$ & Near ($1.5\times$) & 200  \\ \hline
LiH & 12 & 631  & 96  & $-7.5463$ & $-7.8825$ & $0.3362$ & No   & 768  \\ \hline
BeH\textsubscript{2} & 14 & 666  & 112 & $-15.5557$ & $-15.5951$ & $0.0394$ & No ($25\times$) & 896  \\ \hline
H\textsubscript{2}O & 14 & 1086 & 112 & $-74.7921$ & $-75.0124$ & $0.2203$ & No & 896  \\ \hline
\end{tabular}}
\end{table} 

\begin{table}[htbp]
\centering
\renewcommand{\arraystretch}{1.2}
\caption{Serial baseline VQE results (single process, no MPI, matching SPSA hyperparameters with minor ansatz differences noted in Section 4, seed=42). Both the distributed stack and serial baseline use near zero parameter initialization, track the best-physical energy.}
\label{table:serial_baseline}
\begin{tabular}{|l|r|r|r|r|r|}
\hline
\textbf{Molecule} & \textbf{Energy ($E_h$)} & \textbf{FCI ($E_h$)} & \textbf{Error ($E_h$)} & \textbf{Iters} & \textbf{Wall Time (s)} \\ \hline
H\textsubscript{2} & $-1.1258$ & $-1.1373$ & $0.0115$ & 200 & $0.64$  \\ \hline
LiH & $-7.8803$ & $-7.8825$ & $0.0022$ & 576 & $41.62$ \\ \hline
BeH\textsubscript{2} & $-15.5557$ & $-15.5951$ & $0.0394$ & 896 & $173.84$ \\ \hline
H\textsubscript{2}O & $-74.7892$ & $-75.0124$ & $0.2232$ & 896 & $227.23$ \\ \hline
\end{tabular}
\end{table}

In the stack, the benchmark molecule $H_{2}$ achieved an absolute error of 2.4 mHa, approaching the chemical accuracy threshold 
of 1.6 mHa (1 kcal/mol). Notably from these results, the serial baseline achieved a lower error on $LiH$ (2.2 mHa compared to the much higher 336.2 mHa) 
and a comparable error for $BeH_2$ (39.4 mHa for both paths). This indicates the integration of MPI distribution does not necessarily improve the stack's accuracy, as both tested versions of the 
stack use identical near zero initialization and the matching SPSA hyperparameters, with minor ansatz differences noted in Section 4.
 The molecule $BeH_{2}$ obtained 39.4 mHa ($25\times$ the chemical accuracy threshold), indicating that larger molecules require either more iterations or particle-conserving ansatzes to approach chemical accuracy.
Based on the results of the wall-clock comparison (Table~\ref{table:speedup}), it can be concluded that the MPI distributed stack achieves a meaningful speedup, 
but only on larger molecules. The speedups achieved were $1.11\times$ on $BeH_{2}$ (666 Pauli terms) and $1.18\times$ on $H_{2}O$ (1086 terms), where 
the parallelized Pauli term evaluation will outweigh the overhead from MPI coordination. For the smaller benchmark molecules of $H_{2}$, $LiH$, the 
overhead from MPI initialization and \texttt{Allreduce} communication exceeds the benefit of distributing fewer Pauli terms.
For $H_{2}O$, the distributed stack at $P=2$ achieved a throughput of $896 / 193.01 = 4.65$ iter/s, compared to the serial baseline's $896 / 227.23 = 3.94$ iter/s.

\begin{table}[htbp]
\centering
\renewcommand{\arraystretch}{1.2}
\caption{Wall-clock comparison between the serial baseline and distributed stack at $P=2$ MPI ranks. The speedup is calculated by $T_{\text{serial}} / T_{\text{distributed}}$. This distributed stack achieved meaningful speedup, only for molecules that have a sufficiently large number of distributable Pauli terms.}
\label{table:speedup}
\begin{tabular}{|l|r|r|c|r|r|}
\hline
\textbf{Molecule} & \textbf{Serial (s)} & \textbf{Distributed $P$=2 (s)} & \textbf{Speedup} & \textbf{Serial/iter (s)} & \textbf{Dist./iter (s)} \\ \hline
H\textsubscript{2} & $0.64$  & $2.25$  & $0.28\times$ & $0.003$ & $0.011$ \\ \hline
LiH & $41.62$ & $66.54$  & $0.63\times$ & $0.072$ & $0.087$ \\ \hline
BeH\textsubscript{2} & $173.84$ & $155.89$ & $\mathbf{1.11\times}$ & $0.194$ & $0.174$ \\ \hline
H\textsubscript{2}O & $227.23$ & $193.01$ & $\mathbf{1.18\times}$ & $0.254$ & $0.215$ \\ \hline
\end{tabular}
\end{table}

\subsection{Masking Metric}
The masking metric $M = T_{\text{accel}} / T_{\text{comm}}$  was measured once each iteration for all benchmarked molecules. The values are shown in Table~\ref{table:masking}.

\begin{table}[htbp]
\centering
\renewcommand{\arraystretch}{1.2}
\caption{Communication masking metric $M = T_{\text{accel}} / T_{\text{comm}}$ across the execution paths. $M \geq 1$ indicates that the classical computation fully masks communication latency. In the simulator path, $T_{\text{comm}}$ is the MPI \texttt{Allreduce} time, however in the IBM path, $T_{\text{comm}}$ is the QPU RTT.}
\label{table:masking}
\begin{tabular}{|l|c|>{\raggedright\arraybackslash}p{6cm}|}
\hline
\textbf{Execution Path} & \textbf{$M$ Range} & \textbf{Interpretation} \\ \hline
Simulator $P$=1 & $2{,}200$-$9{,}000$  & Compute dominates by $\sim 3$ orders of magnitude  \\ \hline
Simulator $P$=2 & $12$-$5{,}400$  & Compute dominates, higher variance from MPI timing  \\ \hline
IBM QPU & $<1$ & QPU RTT (32-60\,s) dominates, no concurrent classical overlap in the current build \\ \hline
\end{tabular}
\end{table}

All the tested molecules achieved $M \gg 1$ on the simulator path, confirming that both the statevector construction and Pauli term evaluation 
(the acceleration workload) completely mask the MPI communication overhead. This result validates this project's core 
architectural claim, that distributing a classical subroutine across several MPI ranks will not introduce a communication 
bottleneck, so the CPU-QPU bottleneck \parencite{cruise2020practicalquantumcomputingvalue} is not the main limiting factor of the stack's performance when compute dominates.

\subsection{Strong Scaling}
\label{sec:results_scaling}
\begin{table}[!ht]
\centering
\renewcommand{\arraystretch}{1.2}
\begin{minipage}{\textwidth}
\centering
\caption{Per molecule wall-clock time across the MPI rank counts, where all runs executed on one Docker host with shared CPU cores. $P \geq 4$ results represent a lower bound on the multi-node performance, where ranks would have dedicated resources. $H_{2}$ achieved $0.13$ mHa error at $P=4$, within the $1.6$ mHa threshold.}
\label{table:scaling_molecules}
\begin{tabular}{|l|r|r|r|r|}
\hline
\textbf{Molecule} & \textbf{$P$=1 (s)} & \textbf{$P$=2 (s)} & \textbf{$P$=4 (s)} & \textbf{$P$=8 (s)} \\ \hline
H\textsubscript{2} & $1.11$   & $2.14$  & $3.09$   & $3.69$ \\ \hline
LiH & $77.78$  & $67.84$  & $119.05$  & $113.78$  \\ \hline
BeH\textsubscript{2} & $192.23$ & $196.31$ & $469.92$ & $1030.21$ \\ \hline
H\textsubscript{2}O & $253.06$ & $252.79$ & $488.25$ & $1038.50$ \\ \hline
\end{tabular}
\end{minipage}

\vspace{1.2em}

\begin{minipage}{\textwidth}
\centering
\caption{Strong scaling efficiency for molecule $H_{2}O$ (1086 Pauli terms, 896 iterations). Efficiency $E = T_1/ (P \times T_P)$, $P=2$ result demonstrates meaningful parallelism. The degradation occurring at $P \geq 4$ is due to the shared memory contention occurring on a single Docker host.}
\label{table:scaling_efficiency}
\begin{tabular}{|c|r|c|c|c|}
\hline
\textbf{Ranks ($P$)} & \textbf{$T_{\text{total}}$ (s)} & \textbf{Speedup ($T_1/T_P$)} & \textbf{Ideal Speedup} & \textbf{Efficiency} \\ \hline
1 & $253.06$ & $1.00\times$ & $1.00\times$ & $100\%$  \\ \hline
2 & $252.79$ & $1.00\times$ & $2.00\times$ & $50.1\%$ \\ \hline
4 & $488.25$ & $0.52\times$ & $4.00\times$ & $13.0\%$ \\ \hline
8 & $1038.50$ & $0.24\times$ & $8.00\times$ & $3.0\%$  \\ \hline
\end{tabular}
\end{minipage}
\end{table}

Table~\ref{table:scaling_molecules} shows that the CPU-only scaling provides negligible benefit even at $P=2$, with $H_{2}O$ achieving only $1.00\times$ speedup at $50.1\%$ efficiency (Table~\ref{table:scaling_efficiency}). The scaling degrades further at $P=4$ ($0.52\times$) and $P=8$ ($0.24\times$), demonstrating that on CPU-only shared-memory hardware, the MPI distribution overhead exceeds the benefit of distributing Pauli terms. This regression is due to two main factors:   
\begin{enumerate}
    \item \textbf{Statevector redundancy}: Each rank constructs a full $2^{n}$ dimensional statevector (the $O(2^n)$ computational cost) independently and without using parallelization. Only the Pauli term evaluations are distributed.
    \item \textbf{Shared memory contention}: All MPI ranks execute on a single physical host that shares CPU cache, memory bandwidth, and the OS scheduling. This bottleneck would be reduced by the usage of a multi-node cluster with an InfiniBand interconnect.          
\end{enumerate}

The scaling behavior is expected to see significant improvements when handling more complex molecules and by GPU accelerated statevector construction. It should be noted that $H_{2}O$ achieved $0.13$ mHa error at $P=4$ within the strong scaling run data, demonstrating the stack's ability to 
reach the accuracy target in favorable SPSA trajectories, not as a guaranteed result across different rank counts or seeds. 

\subsection{Weak Scaling}
\label{sec:results_weak_scaling}
\begin{table}[htbp]
\centering
\renewcommand{\arraystretch}{1.2}
\caption{Weak scaling results: molecule complexity increases proportionally with rank count to maintain an increasing per-rank workload. All runs executed 10 iterations on a single Docker host with a fixed seed.}
\label{table:weak_scaling}
\begin{tabular}{|c|l|c|c|r|r|r|r|}
\hline
\textbf{$P$} & \textbf{Molecule} & \textbf{Pauli} & \textbf{Terms/Rank} & \textbf{$T_{\text{total}}$ (s)} & \textbf{$T_{\text{iter}}$ (s)} & \textbf{$T_{\text{comm}}$/iter (s)} & \textbf{$M$ (typical)} \\ \hline
1 & H\textsubscript{2} & 15 & 15 & $0.051$  & $0.005$ & $0.000025$ & $198$   \\ \hline
2 & LiH  & 631   & 316 & $0.576$ & $0.058$ & $0.000041$ & $1{,}415$ \\ \hline
4 & BeH\textsubscript{2} & 666 & 167 & $2.411$  & $0.241$ & $0.045$  & $4.3$   \\ \hline
8 & H\textsubscript{2}O & 1086 & 136 & $14.398$ & $1.440$ & $0.256$ & $4.6$   \\ \hline
\end{tabular}
\end{table}

The total wall-clock time grows with problem size, so that the $O(2^n)$ 
statevector construction is the primary computational cost. The MPI communication overhead grows sublinearly, with $T_{\text{comm}}$ per 
iteration increasing from $25\;\mu$s ($P=1$) to $256$\,ms ($P=8$), however, the per term communication cost actually decreases 
from $\sim$17\,$\mu$s/term to $\sim$2.4\,$\mu$s/term. This confirms that \texttt{MPI\_Allreduce} scales efficiently.

The masking metric $M$ drops from $\sim$1{,}400 at $P=2$ down to $\sim$4.5 at $P \geq 4$, reestablishing the shared memory 
contention effect that was observed within the previous strong scaling results. With 4--8 MPI ranks competing for the 4 physical CPU cores, the value of  
$T_{\text{comm}}$ is increased relative to $T_{\text{accel}}$. If this was conducted on dedicated multinode hardware with per rank CPU resources, $M$ 
would remain higher. An important note from the stack's results is how the metric $M > 1$ for all rank counts, confirming the results that even in the worst case, shared memory 
conditions, classical computation still exceeds the communication overhead.

\subsection{IBM Quantum QPU Results}
Table~\ref{table:ibm_results} is the summary of results from the IBM QPU run, while Table~\ref{table:ibm_iterations} 
provides information of each iteration. Over the 10 iterations, the best energy improved from $-0.111\;E_h$ to $-0.171\;E_h$, 
proving a consistent downward trend that confirms the SPSA optimizer configuration is fully functional on real quantum hardware. 
However, the final energy of $-0.171\;E_h$ is far from the expected $H_2$ FCI reference $-1.137\;E_h$ (error of $0.966\;E_h$). 
This was an expected result for several reasons, most importantly that 10 iterations is not enough for SPSA convergence (as the simulator path required 200 
iterations for $H_{2}$). 4096 shot sampling also introduces significant variance, which would not exist in a noiseless statevector
path and T-REx with a resilience level 1 only handles quantum measurement readout errors, but does not address gate errors or decoherence 
on the physical qubits in the IBM runs. The energy trajectory visible in the results of Table~\ref{table:ibm_iterations} 
(where the energy increases at iterations 3,4, and 8) is expected of SPSA operating under shot noise, as each energy estimate 
has a statistical uncertainty of order $1/\sqrt{4096} \approx 0.016$.    

Each iteration's wall-clock time averaged at $42.1$s, with the metric primarily dominated by the QPU RTT (31.9-60.4 s). This finding confirms 
the cloud QPU latency overhead that was predicted in the Methodology, where the QPU RTT exceeded the classical computation time by 
orders of magnitude as when compared to the simulator path time. The stack's architecture handled this unwanted result without causing failure across all 10 iterations.

A final GPU-accelerated QPU experiment was conducted, validating the simultaneous full triple heterogeneous execution path (CPU, GPU, QPU). Over the 10 iterations, this run achieved $-0.163\;E_h$ at $53.2$s/iter,
as compared to CPU-QPU runs ($-0.171\;E_h$ at $42.1$s/iter), with a higher per-iteration time in the smaller molecule $H_{2}$'s workload due to GPU initialization overhead. This result is consistent with the GPU 
overhead results within Section~\ref{sec:gpu_results} for small molecules.


\begin{table}[htbp]
\centering
\renewcommand{\arraystretch}{1.2}
\caption{IBM Quantum QPU results for $H_{2}$ on the \texttt{ibm\_marrakesh} (156 qubit Heron processor), EstimatorV2 with 4096 shots, T-REx measurement error mitigation. The large error compared to FCI is expected with only 10 iterations on the noisy QPU hardware without advanced error mitigation put in place.}
\label{table:ibm_results}
\begin{tabular}{|l|r|r|r|c|r|r|}
\hline
\textbf{Molecule} & \textbf{Energy ($E_h$)} & \textbf{FCI ($E_h$)} & \textbf{Error ($E_h$)} & \textbf{Iters} & \textbf{Wall Time (s)} & \textbf{Time/Iter (s)} \\ \hline
H\textsubscript{2} & $-0.1709$ & $-1.1373$ & $0.9663$ & 10 & $421.0$ & $42.1$ \\ \hline
\end{tabular}
\end{table}

\begin{table}[htbp]
\centering
\renewcommand{\arraystretch}{1.2}
\caption{Per-iteration data from IBM QPU run, with each iteration submitting two PUBs (Primitive Unified Blocs) for both $E(\theta^+)$ and $E(\theta^-)$ gradient estimations. $\Delta E$ is the absolute energy change from the previous iteration. QPU RTT is the round-trip time for each cloud job submission.}
\label{table:ibm_iterations}
\begin{tabular}{|c|r|r|r|l|}
\hline          
\textbf{Iter} & \textbf{Energy ($E_h$)} & \textbf{$\Delta E$ ($E_h$)} & \textbf{QPU RTT (s)} & \textbf{Job ID} \\ \hline
1  & $-0.1108$ & ---      & $31.9$ & \texttt{d6tl3g8v...} \\ \hline
2  & $-0.1180$ & $7.1\mathrm{e}{-3}$  & $33.2$ & \texttt{d6tl3o0v...} \\ \hline
3  & $-0.1161$ & $1.9\mathrm{e}{-3}$  & $32.9$ & \texttt{d6tl40at...} \\ \hline
4  & $-0.1072$ & $8.9\mathrm{e}{-3}$  & $33.4$ & \texttt{d6tl48k6...} \\ \hline
5  & $-0.1328$ & $2.6\mathrm{e}{-2}$  & $48.6$ & \texttt{d6tl4h2f...} \\ \hline
6  & $-0.1347$ & $1.9\mathrm{e}{-3}$  & $41.6$ & \texttt{d6tl4t0v...} \\ \hline
7  & $-0.1621$ & $2.7\mathrm{e}{-2}$  & $60.4$ & \texttt{d6tl57if...} \\ \hline
8  & $-0.1461$ & $1.6\mathrm{e}{-2}$  & $46.5$ & \texttt{d6tl5mk6...} \\ \hline
9  & $-0.1469$ & $8.4\mathrm{e}{-4}$  & $36.9$ & \texttt{d6tl628v...} \\ \hline
10 & $-0.1709$ & $2.4\mathrm{e}{-2}$  & $49.3$ & \texttt{d6tl6bit...} \\ \hline       
\end{tabular}
\end{table}

\subsection{GPU-Accelerated Results}
\label{sec:gpu_results}

Table~\ref{table:gpu_speedup} compares the wall-clock times of the GPU Aer backend against the serial baseline 
across all four benchmark molecules. Unlike the CPU only distributed results (Table~\ref{table:speedup}), 
the GPU accelerated stack achieves meaningful speedup at $P=4$ and continues to improve at $P=8$ instead of providing diminishing returns.

\begin{table}[!ht]
\centering
\renewcommand{\arraystretch}{1.2}
\caption{Wall-clock comparison between the serial baseline and the GPU Aer distributed stack at $P \in \{2,4,8\}$. The GPU backend uses Qiskit Aer with CUDA thrust on an NVIDIA GTX 1650 Mobile. Speedup is calculated as $T_{\text{serial}} / T_{\text{GPU}}$.}
\label{table:gpu_speedup}
\resizebox{\textwidth}{!}{%
\begin{tabular}{|l|r|r|r|r|c|c|c|}
\hline
\textbf{Molecule} & \textbf{Serial (s)} & \textbf{GPU $P$=2 (s)} & \textbf{GPU $P$=4 (s)} & \textbf{GPU $P$=8 (s)} & \textbf{Speedup $P$=2} & \textbf{Speedup $P$=4} & \textbf{Speedup $P$=8} \\ \hline
H\textsubscript{2} & $0.64$  & $0.89$  & $0.89$  & $1.13$ & $0.72\times$ & $0.72\times$ & $0.57\times$ \\ \hline
LiH & $41.62$ & $54.41$ & $44.36$ & $38.74$ & $0.76\times$ & $0.94\times$ & $\mathbf{1.07\times}$ \\ \hline
BeH\textsubscript{2} & $173.84$ & $152.24$ & $148.14$ & $141.16$ & $\mathbf{1.14\times}$ & $\mathbf{1.17\times}$ & $\mathbf{1.23\times}$ \\ \hline
H\textsubscript{2}O & $227.23$ & $187.72$ & $160.72$ & $150.25$ & $\mathbf{1.21\times}$ & $\mathbf{1.41\times}$ & $\mathbf{1.51\times}$ \\ \hline
\end{tabular}}
\end{table}

The critical difference from the CPU only results is the scaling behavior at $P \geq 4$, as the CPU only 
stack regressed down to $0.52\times$ at $P=4$ and $0.24\times$ at $P=8$ for the molecule $H_{2}O$ (Table~\ref{table:scaling_efficiency}), whereas 
the GPU Aer backend achieves an expected speedup with the $1.41\times$ at $P=4$ and $1.51\times$ at $P=8$. This finding confirms our hypothesis 
from Section~\ref{sec:results_scaling}, that the $O(2^n)$ statevector construction was the primary bottleneck and the introduction of 
GPU acceleration of gate operations reduces this cost sufficiently for MPI Pauli term distribution to 
provide continued computational benefit at higher rank counts.

Table~\ref{table:gpu_scaling} presents the GPU Aer strong scaling efficiency for $H_{2}O$.

\begin{table}[!ht]
\centering
\renewcommand{\arraystretch}{1.2}
\caption{Strong scaling efficiency for $H_{2}O$ with GPU Aer backend (1086 Pauli terms, 896 iterations). Compared to the CPU-only efficiency of $50.1\%$ at $P=2$ and $3.0\%$ at $P=8$, the GPU path maintains meaningful efficiency across all rank counts.}
\label{table:gpu_scaling}
\begin{tabular}{|c|r|c|c|}
\hline
\textbf{Ranks ($P$)} & \textbf{$T_{\text{total}}$ (s)} & \textbf{Speedup ($T_1/T_P$)} & \textbf{Efficiency} \\ \hline
1 & $266.90$ & $1.00\times$ & $100\%$  \\ \hline
2 & $187.72$ & $1.42\times$ & $71.1\%$ \\ \hline
4 & $160.72$ & $1.66\times$ & $41.5\%$ \\ \hline
8 & $150.25$ & $1.78\times$ & $22.2\%$ \\ \hline
\end{tabular}
\end{table}

\subsubsection{cuStateVec Backend Comparison}
The cuStateVec backend, where the full statevector resides in GPU memory, was consistently slower than the 
GPU Aer backend on the test hardware (NVIDIA GTX 1650 Mobile, 4 GB GDDR6). For $H_{2}O$ at $P=2$, cuStateVec 
took $183.49$\,s compared to GPU Aer's $182.85$\,s, with the gap widening at $P=4$ ($242.11$\,s vs. $165.76$\,s). 
This result is explained by the hardware mismatch: \textcite{bayraktar2023cuquantum} established 
that statevector simulation is memory-bandwidth-bound, and the GTX 1650's $\sim$128 GB/s GDDR6 bandwidth is 
over $15\times$ lower than the A100's 2,039 GB/s HBM2e. The overhead of GPU kernel launches and data transfers on this consumer hardware exceeds the benefit for circuits of 14 qubits or fewer. On data-center GPUs with HBM, cuStateVec would be expected to outperform the CPU Aer path significantly.

\subsubsection{GPU Weak Scaling}
Table~\ref{table:gpu_weak_scaling} presents the weak scaling results with GPU acceleration compared to the CPU only path.

\begin{table}[htbp]
\centering
\renewcommand{\arraystretch}{1.2}
\caption{Weak scaling comparison: CPU-only vs. GPU Aer. GPU acceleration reduces $H_{2}O$ at $P=8$ from 14.4\,s to 2.4\,s, a $6\times$ improvement.}
\label{table:gpu_weak_scaling}
\begin{tabular}{|c|l|r|r|r|r|}
\hline
\textbf{$P$} & \textbf{Molecule} & \textbf{CPU $T_{\text{total}}$ (s)} & \textbf{GPU $T_{\text{total}}$ (s)} & \textbf{CPU $T_{\text{iter}}$ (s)} & \textbf{GPU $T_{\text{iter}}$ (s)} \\ \hline
1 & H\textsubscript{2}  & $0.051$ & $0.050$ & $0.005$ & $0.005$ \\ \hline
2 & LiH & $0.576$ & $0.719$ & $0.058$ & $0.072$ \\ \hline
4 & BeH\textsubscript{2} & $2.411$ & $1.793$ & $0.241$ & $0.179$ \\ \hline
8 & H\textsubscript{2}O  & $14.398$ & $2.389$ & $1.440$ & $0.239$ \\ \hline
\end{tabular}
\end{table}

Within these results, the GPU weak scaling proves that the per iteration time remains relatively stable across all rank 
counts with a range of $0.005$--$0.239$\,s, compared to the CPU path where $T_{\text{iter}}$ grows from $0.005$\,s to $1.440$\,s. 
The $6\times$ improvement at $P=8$ ($H_{2}O$) confirms GPU acceleration reduces the 
statevector bottleneck that dominated the previous results.

It should be noted for the smaller molecules, the GPU kernel launch overhead on the GTX 1650 Mobile exceeds the 
acceleration benefit for circuits that have 12 or less qubits. For example, the molecule $LiH$ at $P=2$ takes $0.719$\,s on GPU vs. $0.576$\,s on CPU. 
The advantage GPU is only observable at $P \geq 4$ with the larger workloads.

\subsection{Checkpoint Resilience}
The implemented 7 layer diagnostic test has verified the checkpoint resilience. After 5 iterations for 
$H_{2}$, the optimizer state ($\theta \in \mathbb{R}^{24}$) saved to \texttt{checkpoint\_iter\_0005.npy} and 
correctly resumed from iteration 6, using the saved $\theta$ and the energy 
trajectory continued on without error. This stack's checkpoint system does rotate the saved files to keep 
only the 5 most recent, preventing unbounded disk usage.

\begin{table}[htbp]
\centering
\renewcommand{\arraystretch}{1.2}
\caption{Checkpoint resilience validation. Every 5 iterations, VQE loop checkpoints the global $\theta$ state. A simulated failure at iteration 5 will trigger a restart from the saved checkpoint to resume the SPSA schedule without any loss of optimization progress.}
\label{table:resilience}
\begin{tabular}{|l|r|r|r|}
\hline
\textbf{Phase} & \textbf{Iterations} & \textbf{Start Energy ($E_h$)} & \textbf{End Energy ($E_h$)} \\ \hline
Initial run (iters 1--5) & 5 & $-0.963$ & $-1.071$ \\ \hline
Post-restart (iters 6--10) & 5 & $-1.085$ & $-1.109$ \\ \hline
\end{tabular}
\end{table}

The Experiment 6 additional stress tests also passed, with the QPU latency spiking test, 
random delays of 0.5--2.0\,s were injected into MPI odd numbered ranks across 10 VQE iterations for 
the molecule $H_{2}$. All iterations completed without a MPI timeout or deadlock, and all returned 
energies were finite. The testing of drop out recovery, where the iteration 10 checkpoint was intentionally 
deleted after 10 iteration, forcing the stack to fall back to the iteration 5 checkpoint. The 
restarted optimization resumed from the correct SPSA schedule position ($k=6$), and  
energy trajectory continued to descend, which confirms no optimization progress was permanently lost.

\begin{table}[!ht]
\centering
\renewcommand{\arraystretch}{1.2}
\caption{Evaluation of all quantitative success metrics (defined in Section~\ref{sec:methodology}) against the observed results within the experiments.}
\label{table:success_evaluation}
\resizebox{\textwidth}{!}{%
\begin{tabular}{|>{\raggedright\arraybackslash}p{2.8cm}|>{\raggedright\arraybackslash}p{1.3cm}|>{\raggedright\arraybackslash}p{3.0cm}|>{\raggedright\arraybackslash}p{3.5cm}|>{\raggedright\arraybackslash}p{2.5cm}|}
\hline
\textbf{Category} & \textbf{Metric} & \textbf{Threshold} & \textbf{Observed} & \textbf{Status} \\ \hline
System Performance & $T_{\text{total}}$ & 1.5$\times$ speedup & $1.51\times$ (H\textsubscript{2}O, GPU $P$=8) & \textbf{Achieved}  \\ \hline
Latency Masking & $M$ & $M \geq 1$ & $M = 12$--$9{,}000$ (simulator) & \textbf{Achieved} (simulator) \\ \hline
Throughput & iter/s & $>$ serial baseline & $4.65$ iter/s (H\textsubscript{2}O, $P$=2) vs. $3.94$ serial & \textbf{Achieved} \\ \hline
Hardware Efficiency & $E$ & $\geq 70\%$ at $P$=2 & $71.1\%$ (GPU, $P$=2); $50.1\%$ (CPU) & \textbf{Achieved} (GPU) \\ \hline
Scientific Fidelity & $\Delta E$ & $\leq 1.6$ mHa & $2.4$ mHa (H\textsubscript{2}); $0.1$ mHa at $P$=4 & Near / \textbf{Partial} \\ \hline
Resilience & $R_{\text{time}}$ & $< 60$\,s recovery & Checkpoint restart verified & \textbf{Achieved} \\ \hline
\end{tabular}}
\end{table}